%; whizzy paragraph -pdf xpdf -latex ./whizzypdfptex.sh
%; whizzy-paragraph "^\\\\begin{frame}\\|\\\\emtext"
% latex beamer presentation.
% platex, latex-beamer でコンパイルすることを想定。 

%     Tokyo Debian Meeting resources
%     Copyright (C) 2012 Junichi Uekawa

%     This program is free software; you can redistribute it and/or modify
%     it under the terms of the GNU General Public License as published by
%     the Free Software Foundation; either version 2 of the License, or
%     (at your option) any later version.

%     This program is distributed in the hope that it will be useful,
%     but WITHOUT ANY WARRANTY; without even the implied warreanty of
%     MERCHANTABILITY or FITNESS FOR A PARTICULAR PURPOSE.  See the
%     GNU General Public License for more details.
%     You should have received a copy of the GNU General Public License
%     along with this program; if not, write to the Free Software
%     Foundation, Inc., 51 Franklin St, Fifth Floor, Boston, MA  02110-1301 USA

\documentclass[cjk,dvipdfmx,12pt]{beamer}
\usetheme{Tokyo}
\usepackage{monthlypresentation}

%  preview (shell-command (concat "evince " (replace-regexp-in-string "tex$" "pdf"(buffer-file-name)) "&")) 
%  presentation (shell-command (concat "xpdf -fullscreen " (replace-regexp-in-string "tex$" "pdf"(buffer-file-name)) "&"))
%  presentation (shell-command (concat "evince " (replace-regexp-in-string "tex$" "pdf"(buffer-file-name)) "&"))

%http://www.naney.org/diki/dk/hyperref.html
%日本語EUC系環境の時
\AtBeginDvi{\special{pdf:tounicode EUC-UCS2}}
%シフトJIS系環境の時
%\AtBeginDvi{\special{pdf:tounicode 90ms-RKSJ-UCS2}}

\newenvironment{commandlinesmall}%
{\VerbatimEnvironment
  \begin{Sbox}\begin{minipage}{1.0\hsize}\begin{fontsize}{8}{8} \begin{BVerbatim}}%
{\end{BVerbatim}\end{fontsize}\end{minipage}\end{Sbox}
  \setlength{\fboxsep}{8pt}
% start on a new paragraph

\vspace{6pt}% skip before
\fcolorbox{dancerdarkblue}{dancerlightblue}{\TheSbox}

\vspace{6pt}% skip after
}
%end of commandlinesmall

\title{Debian 9 Stretch リリースパーティ \\ 東京エリアDebian勉強会}
\subtitle{第152回 2017年6月度}
\author{杉本 典充  dictoss@live.jp}
\date{2017年6月18日}
\logo{\includegraphics[width=8cm]{image200607/openlogo-light.eps}}

\begin{document}

\begin{frame}
\titlepage{}
\end{frame}

\begin{frame}{Agenda}
 \begin{minipage}[t]{0.45\hsize}
  \begin{itemize}
  \item 最近あったDebian関連のイベント報告
	\begin{itemize}
	\item 第151回東京エリアDebian勉強会
	\end{itemize}
  \item 事前課題発表
  \end{itemize}
 \end{minipage}
 \begin{minipage}[t]{0.45\hsize}
  \begin{itemize}
   \item Debian Trivia Quiz
  \end{itemize}
 \end{minipage}
\end{frame}

\section{イベント報告}
\emtext{イベント報告}

\begin{frame}{第151回東京エリアDebian勉強会}
\begin{itemize}
\item 2017年5月20日に、朝日ネットさんをお借りして開催しました
\item 参加者は6名でした
\item 「Debbugs とのつきあいかた:SOAP 編」 発表者：kenhys さん
\end{itemize} 
\end{frame}

\section{事前課題}
\emtext{事前課題}
{\footnotesize
 \begin{prework}{ yamaneco }
  \begin{enumerate}
  \item $B!J2sEz$J$7!K(B
  \item $B!J2sEz$J$7!K(B
  \end{enumerate}
\end{prework}

\begin{prework}{ Charles Plessy }
  \begin{enumerate}
  \item Reproducible, continuously checked, rock solid more than ever.
  \item I hope there will be good updates for hardware support, because hibernation on my laptop still does not work :(
  \end{enumerate}
\end{prework}

\begin{prework}{ Hiroshi Miura }
  \begin{enumerate}
  \item $B!J2sEz$J$7!K(B
  \item $B!J2sEz$J$7!K(B
  \end{enumerate}
\end{prework}

\begin{prework}{ yy\_y\_ja\_jp }
  \begin{enumerate}
  \item $B!J2sEz$J$7!K(B
  \item $B!J2sEz$J$7!K(B
  \end{enumerate}
\end{prework}

\begin{prework}{ Marc Dequ\`enes (Duck) }
  \begin{enumerate}
  \item $B!J2sEz$J$7!K(B
  \item $B!J2sEz$J$7!K(B
  \end{enumerate}
\end{prework}

\begin{prework}{ Yusaku }
  \begin{enumerate}
  \item $B$*$a$G$H$&$4$6$$$^$9!*(B
  \item $B!J2sEz$J$7!K(B
  \end{enumerate}
\end{prework}

\begin{prework}{ Haruki TSURUMOTO(tSU\_RooT) }
  \begin{enumerate}
  \item $B<!$N%j%j!<%9$G$O$b$C$HB?$/$N(B($BM-MQ$J(B)$B%Q%C%1!<%8$r$$$l$F$_$?$$$G$9(B
  \item reproducible-builds$B$N%5%]!<%H(B
  \end{enumerate}
\end{prework}

\begin{prework}{ MichelDaenzer }
  \begin{enumerate}
  \item $B!J2sEz$J$7!K(B
  \item $B!J2sEz$J$7!K(B
  \end{enumerate}
\end{prework}

\begin{prework}{ akedon }
  \begin{enumerate}
  \item $B$*$a$G$H$&$4$6$$$^$9!#(B
  \item $B!J2sEz$J$7!K(B
  \end{enumerate}
\end{prework}

\begin{prework}{ karosuwindam }
  \begin{enumerate}
  \item $B!J2sEz$J$7!K(B
  \item $B!J2sEz$J$7!K(B
  \end{enumerate}
\end{prework}

\begin{prework}{ $BBg?@M4??(B }
  \begin{enumerate}
  \item DM$B!"(BDD$B$N3'MM$O$*Hh$lMM$G$9!#$*$a$G$H$&$4$6$$$^$9!*(B
  \item Linux-4.9$B$O:G$bBg$-$JJQ99$H8@$o$l$F$$$?$N$G!"%+!<%M%k2s$j$N?75!G=$b;n$7$F$_$?$$$G$9(B
  \end{enumerate}
\end{prework}

\begin{prework}{ koedoyoshida }
  \begin{enumerate}
  \item $B$a$G$?$$!#%j%j!<%9$*Hh$lMM$G$9!#(B
  \item $B$b$&;H$C$F$^$9!#(B
  \end{enumerate}
\end{prework}

\begin{prework}{ NOKUBI Takatsugu }
  \begin{enumerate}
  \item $B$O$d$/(Bwheezy$B%^%7%s>e$2$J$-$c!D(B
  \item $B0BDj$7$F$$$l$P$=$l$GK~B-$G$9(B
  \end{enumerate}
\end{prework}

\begin{prework}{ Roger Shimizu }
  \begin{enumerate}
  \item GNU/screen $B$,BP1~$5$l$?(B Debian Installer $B$r;H$C$F8+$F$/$@$5$$(B
  \item $B!J2sEz$J$7!K(B
  \end{enumerate}
\end{prework}

\begin{prework}{ bringer1092 }
  \begin{enumerate}
  \item $B$*$a$G$H$&$4$6$$$^$9(B
  \item $B!J2sEz$J$7!K(B
  \end{enumerate}
\end{prework}

\begin{prework}{ dictoss }
  \begin{enumerate}
  \item $B?7$7$$(Bdebian$B!"$$$d$a$G$?$$!*(B
  \item $BB?$/$N?M$K;H$C$F$b$i$($k$h$&$J(BOS$B$G$"$C$F$[$7$$$G$9!#(B
  \end{enumerate}
\end{prework}

\begin{prework}{ kenhys }
  \begin{enumerate}
  \item $B!J2sEz$J$7!K(B
  \item $B!J2sEz$J$7!K(B
  \end{enumerate}
\end{prework}

\begin{prework}{ Betsuyaku }
  \begin{enumerate}
  \item ibus$B$5$s$NJQKF$K8MOG$$(B
  \item KVM$B$^$o$j$N?J2=(B
  \end{enumerate}
\end{prework}

\begin{prework}{ ysaito }
  \begin{enumerate}
  \item $B!J2sEz$J$7!K(B
  \item $B!J2sEz$J$7!K(B
  \end{enumerate}
\end{prework}

\begin{prework}{ henrich }
  \begin{enumerate}
  \item $B!J2sEz$J$7!K(B
  \item $B!J2sEz$J$7!K(B
  \end{enumerate}
\end{prework}

}

\subsection{問題}
%; whizzy-master ../debianmeetingresume201311.tex
% $B0J>e$N@_Dj$r$7$F$$$k$?$a!"$3$N%U%!%$%k$G(B M-x whizzytex $B$9$k$H!"(Bwhizzytex$B$,MxMQ$G$-$^$9!#(B
%

\santaku
{$B$D$$$K%j%j!<%9$5$l$?(BDebian 9 (Stretch)$B!#%j%j!<%9%N!<%H$K=q$+$l$F$$$k%j%j!<%9$5$l$?%Q%C%1!<%8?t$O$*$h$=$I$N$/$i$$$G$7$g$&$+!#(B}
{40,000 $B%Q%C%1!<%8$/$i$$(B}
{50,000 $B%Q%C%1!<%8$/$i$$(B}
{60,000 $B%Q%C%1!<%8$/$i$$(B}
{B}
{$B%j%j!<%9%N!<%H$N!V(B2.2. $B%G%#%9%H%j%S%e!<%7%g%s$N:G?7>pJs!W(B(\url{https://www.debian.org/releases/stretch/amd64/release-notes/ch-whats-new.html})$B$N5-:\$K$h$k$H!"(B51687$B%Q%C%1!<%8$rD6$($k$H$N$3$H$G$9!JJdB-!'(Bmain$B$N%P%$%J%j%Q%C%1!<%8$N?t$G$9!K!#(BDD$B!"(BDM$B!"(BPM$B$N$_$J$5$^!":n6H$"$j$,$H$&$4$6$$$^$7$?!#(B}
% $ grep '^Package: ' /var/lib/apt/lists/*\_stretch\_main\_binary-amd64* | wc -l
% 50840

\santaku
{$B$^$@OC$OAa$$$G$9$,!"$3$l$+$i$O<!$N(BDebian 10$B$K8~$+$C$F:n6H$r?J$a$k$3$H$K$J$j$^$9!#(BDebian 10$B$N%3!<%I%M!<%`$O$J$s$G$7$g$&$+!#(B}
{Potato}
{Bullseye}
{Buster}
{C}
{Debian 10$B$N%3!<%I%M!<%`$O(BBuster$B$K$J$j$^$9!#(BDebian$B$N%j%j!<%9>pJs$O!V(B\url{https://wiki.debian.org/DebianReleases}$B!W$K$^$H$^$C$F$$$^$9$N$G$4;2>H$/$@$5$$!#(B}


\section{セミナー}
\emtext{セミナー}

%\section{セミナー1}
%\emtext{Debian Policyドキュメント改良の試み}

%\section{セミナー2}
%\emtext{no title}

\section{Hack time}
\emtext{Hack time}

%\begin{frame}{成果を記入下さい!}
%\end{frame}
  
\section{今後のイベント}
\emtext{今後のイベント}
\begin{frame}{今後のイベント}
  \begin{itemize}
  \item 2017/7/14(金) Debian/Ubuntu ユーザミートアップ in 札幌 （予定）
  \item 2017/7/15(土) OSC 2017 Hokkaidoへ出張参加
    \begin{itemize}
    \item 東京エリアDebian勉強会のメンバが参加
    \item セミナーやります
    \end{itemize}
  \item 2017/7の東京エリアDebian勉強会はお休みとなります
  \item 2017/8/19(土) 東京エリアDebian勉強会（予定）
    \begin{itemize}
    \item 日時、会場、セミナー内容は調整中
    \end{itemize}
  \item 勉強会では発表者を随時募集中です。ぜひともよろしくお願いします。
  \end{itemize}
\end{frame}

\end{document}

;;; Local Variables: ***
;;; outline-regexp: "\\([ 	]*\\\\\\(documentstyle\\|documentclass\\|emtext\\|section\\|begin{frame}\\)\\*?[ 	]*[[{]\\|[]+\\)" ***
;;; End: ***
